% ===== 原文 pp.13--17 =====
% [原文 p.13]

\para{沿一般余维-$(p+1)$ 子流形的插入。} 算符 $U$ 可以沿时空的任意余维-$(p+1)$ 子流形 $\Sigma_{d-p-1}$ 插入。插入之后，该算符记为
\begin{equation}
  U(\Sigma_{d-p-1})\,. \tag{2.38}
\end{equation}
特别地，我们可以选取 $\Sigma_{d-p-1}$ 位于某一特定时刻 $t$ 的空间切片内；此时插入的算符 $U(\Sigma_{d-p-1})$ 可以等同于一个作用在该时刻 $t$ 的 Hilbert 空间上的幺正算符。

% [Tutor Note: p.13 (2.38) 要区分 $U$ 与 $U(\Sigma_{d-p-1})$：$U$ 是“可插入任意余维-(p+1) 子流形”的算符族，括号中给出流形即表示“已插入”。核对办法：取 $p=0$，则余维为 1，(2.41) 退化为 $U(t)U^{-1}(t)=1$，与 §2.1 中 0-形式对称性算符的写法完全一致。]

\para{拓扑性质。} $U$ 是拓扑算符这一事实意味着
\begin{equation}
  U(\Sigma_{d-p-1}) = U(\Sigma'_{d-p-1})\,, \tag{2.39}
\end{equation}
只要 $\Sigma'_{d-p-1}$ 是时空的一个余维-$(p+1)$ 子流形、且可由 $\Sigma_{d-p-1}$ 经拓扑形变得到，亦即只要有
\begin{equation}
  \Sigma'_{d-p-1} - \Sigma_{d-p-1} = \partial\Sigma_{d-p}\,, \tag{2.40}
\end{equation}
其中 $\partial\Sigma_{d-p}$ 是时空的 $(d-p)$ 维子流形 $\Sigma_{d-p}$ 的边界。这一情形可以类比图 1 来图示。

\para{可逆性质。} $U$ 可逆这一事实意味着存在另一个拓扑余维-$(p+1)$ 算符 $U^{-1}$，使得
\begin{equation}
  U(\Sigma_{d-p-1})U^{-1}(\Sigma_{d-p-1}) = 1\,. \tag{2.41}
\end{equation}
也就是说，沿同一个余维-$(p+1)$ 子流形 $\Sigma_{d-p-1}$ 同时插入 $U$ 与 $U^{-1}$，等价于沿 $\Sigma_{d-p-1}$ 不插入任何算符。

\para{$p$-形式对称群。} 一般地，$p$-形式对称性组合起来构成一个 $p$-形式对称群
\begin{equation}
  G^{(p)}\,. \tag{2.42}
\end{equation}
这意味着我们有一族拓扑余维-$(p+1)$ 算符
\begin{equation}
  U_g\,, \qquad \forall g\in G^{(p)}\,, \tag{2.43}
\end{equation}
% [原文 p.14]
由群 $G^{(p)}$ 的元素标号。群乘法描述这些算符的复合
\begin{equation}
  U_g(\Sigma_{d-p-1})U_{g'}(\Sigma_{d-p-1}) = U_{gg'}(\Sigma_{d-p-1})\,, \tag{2.44}
\end{equation}
这通常被记作一条\textbf{融合规则}
\begin{equation}
  U_g\otimes U_{g'} = U_{gg'}\,. \tag{2.45}
\end{equation}

% [Tutor Note: p.14 (2.44)--(2.45) 中算符从左到右的次序对应乘积 $gg'$ 的次序；由于 (2.45) 之后紧接着要证明 $G^{(p)}$ 阿贝尔，两种次序最终给出同一元素。可验证的锚点：取 $G^{(p)}=\Z_N$，把 $g$ 写成 $\ee^{2\pi\ii a/N}$、$g'$ 写成 $\ee^{2\pi\ii b/N}$，则 $\ee^{2\pi\ii(a+b)/N}=\ee^{2\pi\ii(b+a)/N}$，与图 4 的结论自洽。]

\para{高阶形式对称群是阿贝尔的。} 一个重要的性质是：对 $p\ge1$，$p$-形式对称群 $G^{(p)}$ 总是阿贝尔的。这是因为我们可以利用拓扑形变来改变两个余维大于一的拓扑算符的先后次序。见图 4。

% [原文 p.15]
\begin{figure}[htbp]
\centering
\includegraphics[width=0.5\textwidth]{fig4.png}
\figtitle{4}{如图中所示，对于 $p\ge1$，可以把与一个 $p$-形式对称性相联系的拓扑算符绕过与另一个 $p$-形式对称性相联系的拓扑算符。这改变这两个拓扑算符的次序，意味着融合后所得的算符相同，$U_{gg'}=U_{g'g}$，从而 $p\ge1$-形式对称群必须是阿贝尔的。}
\end{figure}

% [Tutor Note: p.14 这一论证的关键前提是“余维大于一”：余维-1 算符要交换次序必须在空间中切开一个洞、从而改变插入所在的切片，所以 0-形式对称群可以非阿贝尔。可验证的反例：$p=0$ 时取置换群 $S_3$（例如三个同种标量场上的置换对称性），它非阿贝尔，而 (2.45) 的融合规则照样成立。]

\para{特例：$\U(1)$ $p$-形式对称性。} 把上面讨论的例 2.1 加以推广，$p$-形式对称性的如下特例
\begin{equation}
  G^{(p)} = U(1) \tag{2.46}
\end{equation}
能干净地展示上述所有概念。这样的对称性联系于 Noether 流的高阶形式推广，即一个局域算符 $j^{d-p-1}(x)$，它是时空上的一个 $(d-p-1)$-形式。连续性方程的类比物是：这个 $(d-p-1)$-形式是闭的
\begin{equation}
  \dd j^{d-p-1} = 0\,. \tag{2.47}
\end{equation}
利用这一点，我们可以把构成 $p$-形式对称性 (2.46) 的拓扑余维-$(p+1)$ 算符写成
\begin{equation}
  U_g(\Sigma_{d-p-1}) = \exp\left(\ii\alpha\int_{\Sigma_{d-p-1}} j^{d-p-1}\right)\,, \qquad g = \ee^{\ii\alpha}\in U(1)\,. \tag{2.48}
\end{equation}
利用 (2.47)、并重复 (2.33) 中给出的同一论证，容易证明这些算符是拓扑的。

% [Tutor Note: p.14 (2.47) 中 $j^{d-p-1}$ 的上标是微分形式的次数，不是幂次（原文一贯如此书写，对比 (2.35) 的 $\partial_\mu j^\mu$）。因此 (2.47) 左边是一个 $(d-p)$-形式；$p=0$ 时它退化为 $(d-1)$-形式的守恒流方程，正是 §2.1 用过的对偶流 $\dd j_{d-1}=0$。]

实现 $\U(1)$ 高阶形式对称性的一个具体例子如下。

\begin{examplebox}{例 2.2：Maxwell 理论中的高阶形式对称性}
考虑一个规范群为 $\U(1)$、且不含物质场的 $d$ 维规范理论。这是通常 Maxwell 理论的 $d$ 维推广；代回 $d=4$ 即得到通常的 Maxwell 理论。

设 $A$ 为 $\U(1)$ 规范场，它是时空上的一个 1-形式。其场强为
\begin{equation}
  F = \dd A\,, \tag{2.49}
\end{equation}
% [原文 p.16]
它是时空上的一个 2-形式。由于外微分的平方为零，这个 2-形式是闭的
\begin{equation}
  \dd F = 0\,, \tag{2.50}
\end{equation}
这也称为 Bianchi 恒等式。把 $F$ 作为 Noether 流，我们推断存在一个 $(d-3)$-形式对称性
\begin{equation}
  G^{(d-3)} = U(1)\,, \tag{2.51}
\end{equation}
其拓扑面算符为
\begin{equation}
  U_g^{(m)}(\Sigma_2) = \exp\left(\ii\alpha\int_{\Sigma_2} F\right)\,, \qquad g = \ee^{\ii\alpha}\in U(1)\,. \tag{2.52}
\end{equation}
$U_g^{(m)}$ 中的上标表示：这个 $\U(1)$ $(d-3)$-形式对称性就是所研究的纯 $\U(1)$ 规范理论的\textbf{磁高阶形式对称性}。

此外，由运动方程我们有
\begin{equation}
  \dd\star F = 0\,, \tag{2.53}
\end{equation}
其中 $\star F$ 是一个 $(d-2)$-形式，即场强 $F$ 的 Hodge 对偶。把 $\star F$ 作为 Noether 流，我们推断存在一个 1-形式对称性
\begin{equation}
  G^{(1)} = U(1)\,, \tag{2.54}
\end{equation}
其拓扑余维-2 算符为
\begin{equation}
  U_g^{(e)}(\Sigma_{d-2}) = \exp\left(\ii\alpha\int_{\Sigma_{d-2}} \star F\right)\,, \qquad g = \ee^{\ii\alpha}\in U(1)\,. \tag{2.55}
\end{equation}
我们把上标写在 $U_g^{(e)}$ 中，因为这个 $\U(1)$ 1-形式对称性就是所研究的纯 $\U(1)$ 规范理论的\textbf{电 1-形式对称性}。正如我们将在例 2.3 中看到的，这实质上就是 Gauss 定律。
\end{examplebox}

% [Tutor Note: p.16 (2.53) 的次数核对：$F$ 是 2-形式，Hodge 星把 $k$-形式映到 $(d-k)$-形式，所以 $\star F$ 是 $(d-2)$-形式，$\dd\star F=0$ 是 $(d-1)$-形式方程；它作为 Noether 流给出 $d-p-1=d-2$ 即 $p=1$，与 (2.54) 的“1-形式对称性”吻合。数值锚点：$d=4$ 时 $\dd\star F=0$ 就是 Maxwell 方程 $\partial_\mu F^{\mu\nu}=0$。]

% [Tutor Note: 例 2.2 方框在原文跨 pp.14--16 断开，此处按 STYLE\_GUIDE §7 合并为一个 examplebox，编号 2.2 与内容均未拆分。]

\subsubsection{$p$-形式对称性的作用}

\para{对局域算符的作用。} 如我们上面讨论的，0-形式对称性可以作用在一个 0 维的、点状的局域算符上。这一作用可以理解为：把与 0-形式对称性相联系的余维-1 拓扑算符 $U_g(\Sigma_{d-1})$ 移过该局域算符的结果。

把这一现象加以推广，首先要指出：$p\ge1$-形式对称性不作用于任何局域算符，因为我们可以沿一个与点 $x$ 不相交的余维-$p$ 子流形 $\Sigma_{d-p}$，把相应的余维-$(p+1)$ 拓扑算符 $U_g(\Sigma_{d-p-1})$ 移过局域算符 $\cO(x)$。

% [Tutor Note: p.16 这段省略了一个前提：能选到与 $x$ 不相交的 $\Sigma_{d-p}$ 要求 $d-p\ge1$，即 $p\le d-1$；$p=d$ 时余维为 $d+1$，已超出时空维数，该“对称性”不存在。可验证：$d=4$ 的 $\U(1)$ 规范理论中电/磁高阶形式对称性的阶数 $p=1$ 与 $p=d-3=1$ 都满足此界。]

% [原文 p.17]
\para{对延展算符的作用。} 现在让我们研究高阶形式对称性能否作用在定义于时空的 $q\ge1$ 维子流形 $M_q$ 上的延展算符。

事实上，按上述论证的一个版本，若 $q<p$，则 $p$-形式对称性不能作用于 $q$ 维算符，因为此时我们可以选取 $\Sigma_{d-p}$ 使之不与 $M_q$ 相交。因此，$p$-形式对称性只能作用于维数为
\begin{equation}
  q\ge p \tag{2.56}
\end{equation}
的算符。最简单的研究对象是 $q=p$ 的情形，此时我们有如下陈述：

\begin{statementbox}{陈述 2.2：$p$ 维算符在 $p$-形式对称性下的荷}
$p$ 维算符按 $p$-形式对称群 $G^{(p)}$ 的表示变换。
\end{statementbox}

让我们更详细地讨论这一作用。考虑一个沿时空的 $p$ 维子流形 $M_p$ 放置的 $p\ge1$ 维延展算符 $\cO(M_p)$。我们可以假定 $\cO(M_p)$ 是一个\emph{不可约} $p$ 维算符，即在任意一点 $x\in M_p$ 处能插入的拓扑局域算符只有单位局域算符的倍数\footnote{这意味着 $p$ 维算符 $\cO(M_p)$ 不能写成另外两个 $p$ 维算符之和。这类算符也称为\emph{单}（simple）$p$ 维算符。}。现在注意到，把 $U_g(\Sigma_{d-p-1})$ 形变着穿过 $\cO(M_p)$，会在 $M_p$ 与 $\Sigma_{d-p}$ 的交点 $x$ 处留下一个拓扑局域算符 $\cO(x)$：
\begin{equation}
  U_g(\Sigma_{d-p-1})\cO(M_p) = \cO(x)\cO(M_p)U_g(\Sigma'_{d-p-1})\,. \tag{2.57}
\end{equation}
见图 5。由于沿 $M_p$ 唯一可能的拓扑局域算符都是单位算符的倍数，我们可以在上式里把 $\cO(x)$ 换成一个非零数
\begin{equation}
  U_g(\Sigma_{d-p-1})\cO(M_p) = \phi(g)\times\cO(M_p)U_g(\Sigma'_{d-p-1})\,, \qquad \phi(g)\in\C^\times = \C-\{0\}\,. \tag{2.58}
\end{equation}
由于融合规则 (2.44)，这些数 $\phi(g)$ 必须满足
\begin{equation}
  \phi(g)\phi(g') = \phi(gg')\,, \tag{2.59}
\end{equation}
也就是说，这些数 $\phi(g)$ 给出 $p$-形式对称群 $G^{(p)}$ 的一个一维表示；特别地，这些数 $\phi(g)$ 必须是相位因子
\begin{equation}
  \phi(g)\in U(1)\subset\C^\times\,. \tag{2.60}
\end{equation}
由于 $G^{(p)}$ 是阿贝尔群，它的一维表示恰好就是它的不可约表示。

因此我们得知：一个不可约 $p$ 维算符 $\cO(M_p)$ 按 $p$-形式对称群 $G^{(p)}$ 的不可约表示变换，这就论证了陈述 2.2。

% [Tutor Note: p.17 从 (2.58) 到 (2.60) 的收紧（$\C^\times\to \U(1)$）原文只写了结论，实际用到了“$G^{(p)}$ 紧”这一条件：紧群的有限维表示可取为幺正的，故 $\phi(g)$ 可化为相位。可验证的反例：若取非紧的 $G^{(p)}=\Z$，则 $\phi(g)=2^{g}$ 满足 (2.59)（$2^{g}2^{g'}=2^{g+g'}$）却不在 $\U(1)$ 中，说明紧性假设不可省。]
